\section*{第 \ref{ch:b2:reaction-enthalpy} 章——反应焓}

\begin{solution}{exo:b2:reaction-enthalpy:1}
$\Delta_r H^\circ = 3(-393.51) + 4(-285.83) - (-104.7) - 5(0) =
\qty{-2219.2}{kJ/mol}$，为负，属放热反应。（实测燃烧
焓与此相同，为 \qty{-2219.2}{kJ/mol}。）
\end{solution}

\begin{solution}{exo:b2:reaction-enthalpy:2}
甲烷：$890.3/16.04 = \qty{55.5}{kJ/g}$，即
每千克 \qty{55.5}{MJ}。氢气：$285.83/2.016 = \qty{141.8}{kJ/g}$，
即 \qty{141.8}{MJ/kg}，每千克约多 2.6 倍（按每升计则
少得多，因气体很轻）。
\end{solution}

\begin{solution}{exo:b2:reaction-enthalpy:3}
$\Delta\nu_{\text{气体}} = 6 - 7.5 = -1.5$，故 $\Delta_r H = \Delta_r U +
\Delta\nu_{\text{气体}}RT = -3264 - 1.5 \times 8.314 \times 298.15 \times
10^{-3} = \qty{-3267.7}{kJ/mol}$，与由生成焓算出的
\qty{-3267.6}{kJ/mol} 相符。
\end{solution}

\begin{solution}{exo:b2:reaction-enthalpy:4}
目标反应等于第一个反应减去第二个反应的两倍：
$\Delta_r H^\circ = -393.5 - 2(-283.0) = \qty{+172.5}{kJ/mol}$。该反应为
吸热：热焦炭在将二氧化碳转化为一氧化
碳时降温。
\end{solution}

\begin{solution}{exo:b2:reaction-enthalpy:5}
断开的键为 C--H 和 Cl--Cl，$416 + 243 = \qty{659}{kJ/mol}$；形成的
键为 C--Cl 和 H--Cl，$350 + 432 = \qty{782}{kJ/mol}$：估算得
\qty{-123}{kJ/mol}。由生成焓得：$-83.7 - 92.3 + 74.9 =
\qty{-101.1}{kJ/mol}$。\qty{22}{kJ/mol} 的差距，来自
平均键焓（对多种分子取平均），以及甲烷中断裂的
C--H 键，比四个键的平均值更强。
\end{solution}

\begin{solution}{exo:b2:reaction-enthalpy:6}
在 \qty{298}{K} 时：$-241.83 + 285.83 = \qty{44.00}{kJ/mol}$。用基尔霍夫定律，取
$\Delta C_p = 33.59 - 75.35 = \qty{-41.76}{J/(K.mol)}$:
$44.00 - 0.04176 \times 101.85 = \qty{39.75}{kJ/mol}$，即 \qty{400}{K} 时的值。
查表可得 $-242.85 + 282.59 = \qty{39.74}{kJ/mol}$：恒定热容
的近似在此区间内效果极佳。
\end{solution}

\begin{solution}{exo:b2:reaction-enthalpy:7}
电离：$4.341 \times 96.485 = \qty{418.8}{kJ/mol}$；电子亲和：
$-3.613 \times 96.485 = \qty{-348.6}{kJ/mol}$。晶格焓 $= 436.7 +
89.0 + 121.3 + 418.8 - 348.6 = \qty{717}{kJ/mol}$，小于
氯化钠的 \qty{787}{kJ/mol}：\ce{K+} 比 \ce{Na+} 大，离子
间距更远，相互吸引更弱。
\end{solution}

\begin{solution}{exo:b2:reaction-enthalpy:8}
量热计的热容：$C = 2.00/0.418 = \qty{4.785}{kJ/K}$。
放出的热量：$4.785 \times 0.583 = \qty{2.789}{kJ}$，对应 $\xi =
\qty{0.0500}{mol}$（酸为限量试剂，碱略过量）：
$\Delta_r H = -2.789/0.0500 = \qty{-55.8}{kJ/mol}$.
\end{solution}

\begin{solution}{exo:b2:reaction-enthalpy:9}
\ce{C8H18(l) + 25/2O2(g) -> 8CO2(g) + 9H2O(l)}: $\Delta_r H^\circ =
8(-393.51) + 9(-285.83) + 250.3 = \qty{-5470.3}{kJ/mol}$；取 $M =
\qty{114.2}{g/mol}$，得 \qty{47.9}{MJ/kg}。二氧化碳：$8 \times 44.01 =
\qty{352.1}{g}$ 对应 \qty{5.470}{MJ}，即 \qty{64.4}{g/MJ}，与
甲烷的 $44.01/0.8903 = \qty{49.4}{g/MJ}$ 相比，甲烷每个
碳原子所含的氢更多，燃烧氢放出热量时不产生二氧化碳。
\end{solution}

\begin{solution}{exo:b2:reaction-enthalpy:10}
\ce{CH4 + 2O2 -> CO2 + 2H2O(g)} 的生成物：$C_p = 37.13 + 2 \times 33.59 =
\qty{104.3}{J/K}$; $q = \qty{802.3}{kJ}$; $T_f = 298 + 802\,300/104.3 \approx
\qty{7990}{K}$，几乎为空气中的三倍，因为没有氮气
来分担热量。真实的氧气火焰要冷得多：远在 \qty{7990}{K} 之前，二氧化碳和水
解离为一氧化碳、氢气、氧气和原子，这些都是吸热
反应，吸收了大部分热量；热容
随 $T$ 增大。
\end{solution}

\begin{solution}{exo:b2:reaction-enthalpy:11}
$\Delta_r H^\circ(700) = \Delta_r H^\circ(298) + \int_{298}^{700}(\Delta_r a +
\Delta_r b\,T)\,\dd T = -91.88 + [-60.5 \times 401.85 + 0.0270 \times
(700^2 - 298.15^2)] \times 10^{-3} = -91.88 - 24.31 + 10.83 =
\qty{-105.4}{kJ/mol}$，与表中值
（\qty{-105.2}{kJ/mol}）相差不超过 \qty{0.2}{kJ/mol}，而恒定 $C_p$ 近似的误差为
\qty{4.5}{kJ/mol}。
\end{solution}

\begin{solution}{exo:b2:reaction-enthalpy:12}
$\Delta_r H^\circ = 2 \times 90.29 = \qty{+180.6}{kJ/mol}$（吸热）。
$\Delta_r C_p^\circ = 2(29.85) - 29.12 - 29.38 = \qty{1.2}{J/(K.mol)}$，故
在 \qty{298}{K} 与 \qty{2000}{K} 之间，反应焓的
变化约为 $1.2 \times 1702 = \qty{2}{kJ/mol}$，仅占百分之一。吸热
反应在高温下更有利（见\cref{ch:b2:equilibrium-shifts}）：
冷空气几乎不生成一氧化氮，最热的火焰生成得
最多，因此燃烧温度是抑制这一
污染物的一个杠杆。
\end{solution}

\begin{solution}{pb:b2:reaction-enthalpy:1}
\textbf{1.} \ce{C4H10(g) + 13/2O2(g) -> 4CO2(g) + 5H2O(l)}.
\textbf{2.} $\Delta_r H^\circ = 4(-393.51) + 5(-285.83) + 125.6 =
\qty{-2877.6}{kJ/mol}$.
\textbf{3.} 若生成物中的水为蒸气，$4(-393.51) + 5(-241.83) + 125.6 =
\qty{-2657.6}{kJ/mol}$；其差值 \qty{220.0}{kJ} 即为
5 摩尔水的汽化焓（每摩尔 \qty{44.0}{kJ/mol}）。
\textbf{4.} $2877.6/58.12 = \qty{49.5}{kJ/g}$ 与 $2657.6/58.12 =
\qty{45.7}{kJ/g}$.
\textbf{5.} 第二个。水以蒸气形式逸出，其
液化热随之损失。
\textbf{6.} $230 \times 45.7 = \qty{1.05e4}{kJ}$，约 \qty{10.5}{MJ}。
\textbf{7.} 在定容条件下，接收的热量为 $\Delta U$（无
体积功）：氧弹量热计测得的是 $\Delta_r U$。
\textbf{8.} $\Delta\nu_{\text{气体}} = 4 - 1 - 6.5 = -3.5$（水
为液态）。
\textbf{9.} $\Delta_r U = \Delta_r H - \Delta\nu_{\text{气体}}RT = -2877.6 +
3.5 \times 2.479 = \qty{-2868.9}{kJ/mol}$.
\textbf{10.} $\xi = 0.500/58.12 = \qty{8.60e-3}{mol}$；热量 $8.60 \times
10^{-3} \times 2868.9 = \qty{24.68}{kJ}$; $\Delta T = 24.68/10.00 =
\qty{2.468}{K}$.
\textbf{11.} 热量 $10.00 \times 2.461 = \qty{24.61}{kJ}$：$\Delta_r U =
-24.61/(8.603 \times 10^{-3}) = \qty{-2860.6}{kJ/mol}$，而 $\Delta_r H =
-2860.6 - 8.7 = \qty{-2869.3}{kJ/mol}$，其绝对值比表中值低
\qty{0.29}{\percent}：热量散失于量热计，或燃烧
略不完全。
\textbf{12.} $n = 1000/18.02 = \qty{55.5}{mol}$; $Q = 55.5 \times 75.35
\times 80 = \qty{334.5}{kJ}$.
\textbf{13.} 放出热量 $16.0 \times 45.72 = \qty{731.5}{kJ}$；效率
$334.5/731.5 = 0.46$.
\textbf{14.} 热量传给从锅边逸出的高温燃烧气体、
被火焰加热的空气、锅壁，以及辐射。
\textbf{15.} $230 \times 45.72 \times 0.457/334.5 = \qty{14.4}{L}$.
\textbf{16.} \qty{6.5}{mol} \ce{O2} 同时带入 $6.5 \times 78.08/20.95 =
\qty{24.23}{mol}$ 的 \ce{N2} 与 $6.5 \times 0.93/20.95 = \qty{0.289}{mol}$
的 \ce{Ar}。
\textbf{17.} \qty{4}{mol} \ce{CO2}, \qty{5}{mol} \ce{H2O(g)},
\qty{24.23}{mol} \ce{N2}, \qty{0.289}{mol} \ce{Ar}.
\textbf{18.} 火焰为绝热等压过程，故 $\Delta H = 0$。沿一条
途径：先在 \qty{298}{K} 下发生反应（$\Delta H_1 =
\qty{-2657.6}{kJ}$），再对这些产物加热，$\Delta H_2 =
-\Delta H_1 = \qty{+2657.6}{kJ}$：热量全部用于加热气体。
\textbf{19.} $C_p = 4(37.13) + 5(33.59) + 24.23(29.12) + 0.289(20.79) =
\qty{1028}{J/K}$.
\textbf{20.} $T_f = 298 + 2\,657\,600/1028 = \qty{2883}{K}$.
\textbf{21.} $T_f = 2300 + 100 \times (2657.6 - 2515.7)/(2655.7 - 2515.7) =
\qty{2401}{K}$.
\textbf{22.} 气体的热容随温度升高（在
\qty{298}{K} 至 \qty{2400}{K} 之间约增加四分之一）：室温下的
数值低估了加热气体所需的热量。
\textbf{23.} 高于 \qty{2000}{K} 时，二氧化碳和水部分解离
（吸热），火焰向外辐射，且与空气的混合永远不可能完美。
\textbf{24.} 丁烷在空气中的绝热火焰温度为
$\boldsymbol{T_{\text{ad}} \approx \qty{2.40e3}{K}}$.
\end{solution}
